\section{Discussion}
The combined evidence supports a disciplined distinction between improving a component and establishing a general solution. The fixed contextual combination does not beat its matching recency control on the selected development blocks. Temporally validated combinations improve the external collection mean point error against that control, while AutoTheta remains stronger on monthly point error and both interval-score comparisons. These outcomes suggest that candidate selection and regularization deserve attention, but they do not isolate the context kernel as the reason for the gain. A faithful comparison with established feature-dependent and arbitrated combination methods remains necessary before making a broader comparative claim.

The browser experiments illustrate a different improvement route. The validation-selected ridge candidate slightly worsens independent point error, so the served point graph remains unchanged. Prefix-stratified residual calibration then improves interval score on an independently generated confirmation collection using that same graph. This gain addresses the uncertainty procedure rather than the point predictor. Its value lies in the paired comparison, fixed evaluation protocol and retained subgroup failures. A stronger point predictor could still be useful, but this experiment supplies no basis for saying that adding regularization automatically creates one.

The daily corpus result similarly separates descriptive association from useful model selection. Features can have reproducible associations with a within-roster error target while a discovery-trained selector performs worse than a simple automatic method on confirmation series. The association target is itself an optimistic oracle over a finite roster, and it can depend on the normalization used by the loss. A feature plot should therefore help formulate a hypothesis rather than trigger an automatic operational choice. The workbench retains both the association table and the unsuccessful selection result so that readers can inspect this distinction.

For the intended research and teaching use, the practical recommendation is to begin with the source, information cutoff and score definition before interpreting a displayed winner. A user should compare point and interval behavior together, inspect individual series as well as aggregate means, and verify whether a displayed result is a fresh fit or an archived replay. Subsequent development can test better-supported alternatives on prospectively specified authentic data and matching resource budgets. The present evidence supports those next experiments and an inspectable evaluation instrument; it does not identify a universally preferred forecast method or an accepted business policy.

\section{Limitations and threats to validity}
The first limitation is sampling. Fifteen deliberately selected cases demonstrate behavior and permit detailed source reconstruction, but they do not estimate accuracy for an identified deployment population. Eleven are generated mechanisms and four are real excerpts. A case with a known generator can reveal leakage, phase errors or an unstable scale; its controlled simplicity can also exclude the missingness, changing definitions and interventions encountered in practice. The larger external studies improve inventory accounting for their specified collections without turning those collections into a random sample of future customer demand.

The second limitation is method comparability. The twenty-two configurations have different context caps, training budgets and historical origins. A direct adaptation may omit an originating paper's architectural components. A maintained package can change default behavior between versions, while a pretrained model can contain related data in an undocumented training corpus. The report addresses implementation identity and source-bound execution, but does not establish that all model families received equally effective tuning or equal computational resources. A common target block is necessary for the stated comparison and insufficient for a general architectural conclusion.

The third limitation concerns statistical interpretation. Origins within a series can be dependent; different series can share a latent source after exact deduplication. The paired-series bootstrap for the external combination study resamples series and is exploratory, with no multiple-comparison adjustment. The disjoint-series feature study controls its declared feature-test family but cannot convert association into causation. The browser confirmation tests generated histories under a fixed mixture of mechanisms. Its marginal coverage and interval-score improvement do not establish authentic-data coverage, simultaneous path coverage or validity within every conditioning group. A remaining trend-family undercoverage is materially relevant to interpretation.

The fourth limitation is engineering and reproducibility scope. Historical streaming hardware details and the original controller-gain rationale were not recovered. Reconstructed arrays can reproduce numerical trajectories without restoring acquisition history or a normalized latency experiment. The internal Python distribution has not yet completed the required reusable-library boundary work, and the existing lesson executions do not establish compliance with the newer teaching-content requirements. Complete observable-unit review and independent reader review remain open. The present document is a working revision, so it is not evidence of a newly accepted publication or completed application release.

